% Separate post-CP17 refinement. B50 remains unchanged during its audit.
\subsection{The exact finite-character list after primitive-norm normalization}

\begin{lemma}[Primitive norm fibres and Frobenius orbits]
\label{lem:B-RD51-orbits}
\status{Proved; complete root and obstruction/crystal proof and exact-code audits passed.}
Let $p\geq5$, and put $m=(p^2-1)/d$, where $d$ is an odd
divisor of $p+1$. Write
\[
                      q=(p+1)/d,\qquad m=(p-1)q.
                                                        \tag{B-RD51.1}
\]
For any fixed primitive $\eta\in\mu_{p-1}$, the fibre
\[
             \{\lambda\in\mu_m:\lambda^{p+1}=\eta\}
                                                        \tag{B-RD51.2}
\]
has $q$ elements and is the disjoint union of $q/2$ orbits
of the involution $\lambda\mapsto\lambda^p$.
Consequently two distinct rank-two tame characters with that
primitive norm and joint faithful order $m$ exist as abstract
characters if and only if $q\geq4$.

In the normalization of Corollary~\ref{cor:B-RD50-remaining},
any unresolved $p^2$ comparison must therefore satisfy
\[
 m=(p-1)q,\qquad q\mid p+1,\qquad
 v_2(q)=v_2(p+1),\qquad q\geq4.                         \tag{B-RD51.3}
\]
The existence assertion in the first paragraph is about finite
characters only; it does not construct Barsotti--Tate groups
or a torsion congruence.
\end{lemma}
\begin{proof}
Since $d$ is odd and divides $p+1$, it is relatively prime to
$p-1$. Hence
\[
             \gcd(m,p+1)=q,\qquad
             |\operatorname{im}(\lambda\mapsto\lambda^{p+1})|
                      =m/q=p-1.
\]
This proves the fibre cardinality. Frobenius preserves the
fibre and has square equal to the identity there, since
$m\mid p^2-1$. A fixed point would belong to $\mu_{p-1}$;
its norm would then be $\lambda^2$, which cannot be primitive
in the even-order group $\mu_{p-1}$. Thus every orbit has
exactly two elements. These orbits are precisely the two
eigenvalues of the rank-two tame characters under discussion.

If $q=2$, the entire norm fibre is a single orbit and no two
different characters occur. Conversely, suppose $q\geq4$.
There is a primitive $m$th root in any specified primitive-norm
fibre. To see this without assuming a preferred generator,
the reduction map $(\mathbf Z/m\mathbf Z)^\times\to
(\mathbf Z/(p-1)\mathbf Z)^\times$ is surjective: lift a unit
modulo $p-1$, and avoid the finitely many additional primes
dividing $m$ by the Chinese remainder theorem. A primitive
$m$th root can therefore be chosen with norm $\eta$.
Choose an element of the fibre outside its Frobenius orbit.
The two resulting characters are distinct, and their joint
order is $m$ because the first one already has order $m$.
Equation (B-RD51.3) now follows from the normalized necessary
list in Corollary~\ref{cor:B-RD50-remaining}.
\end{proof}

\begin{corollary}[Two finite-character patterns in the first remaining prime]
\label{cor:B-RD51-p5}
\status{Proved; complete root and obstruction/crystal proof and exact-code audits passed.}
For the unresolved test $\nu_{\min}(5,4,2)\in\{1,2\}$,
any failure at $25$-torsion has, after the common normalization
above, joint faithful tame degree $24$.
Choose a primitive $24$th root $\zeta$ with
$\zeta^6=\eta$, the prescribed primitive norm. The only
individual character types are
\[
 \begin{array}{c|c|c}
 \text{type}&\text{eigenvalue exponents in }\zeta&
              \text{faithful order}\\ \hline
 A&\{1,5\}&24\\
 B&\{9,21\}&8\\
 C&\{13,17\}&24.
 \end{array}                                             \tag{B-RD51.4}
\]
Up to interchange and simultaneous quadratic twist, there
are only two distinct-type patterns: $\{A,B\}$ and $\{A,C\}$.
The first has individual orders $24,8$; the second consists
of the two opposite quadratic twists of a faithful order-$24$
character. This is an explicit remaining character list, not
an assertion that either pattern is realized by a pair of groups
with isomorphic $25$-torsion.
\end{corollary}
\begin{proof}
The odd divisors of $6$ are $1,3$. The latter gives $q=2$,
which Lemma~\ref{lem:B-RD51-orbits} excludes; hence $m=24$.
The norm fibre consists of the six exponents congruent to
$1$ modulo $4$. Multiplication by $5$ modulo $24$ gives
exactly the three pairs displayed in (B-RD51.4).
A quadratic twist adds $12$ to both exponents, interchanging
$A,C$ and fixing $B$. Therefore the pairs $\{A,B\}$ and
$\{B,C\}$ are equivalent under simultaneous quadratic twist,
while $\{A,C\}$ is the other possibility.
\end{proof}

\begin{remark}[Retained finite calculation and scope]
\label{rem:B-RD51-scope}
\status{Proved under stated hypotheses.}
The exact program
\path{code/reduction_descent/verify_tame_primitive_norm_orbits.py}
enumerates the primitive-norm fibres, their Frobenius orbits,
and the joint orders of distinct pairs for the listed small
primes. It also reproduces (B-RD51.4) and its simultaneous
quadratic-twist orbits. The general assertion is proved above;
the finite calculation is a retained check of the character
arithmetic and makes no group-existence inference. It checked
1,394 primitive-norm fibres and 292,614 distinct-character pairs
for the primes from $5$ through $101$. The author run took
0.406761682999786 seconds; independent root and
obstruction/crystal runs took 0.31523789799393853 and
0.3501194319978822 seconds. Both independent runs preserved
the author output.
\end{remark}
